% =============================================================================
% modules/api.tex — Craftalism API module documentation
% =============================================================================

\modulelang{Java 17 · Spring Boot 3.5 · PostgreSQL}

\chapter{Craftalism API}
\label{ch:api}

\textcolor{CraftGray}{\textit{Core REST API — player registration, balance management, and transaction records}}

\bigskip

% =============================================================================
\section{Overview}
% =============================================================================

The Craftalism API is the central data service for the economy platform. It exposes a typed REST interface under \code{/api} that is consumed by both the Minecraft plugin and the Dashboard. All state for players, balances, and transactions is owned and persisted by this service. Write operations are protected by JWT scope enforcement; the API trusts tokens issued by the Craftalism Authorization Server.

% =============================================================================
\section{Responsibilities}
% =============================================================================

\begin{itemize}
  \item Registers new players and resolves existing players by UUID or display name.
  \item Creates, reads, and overwrites balance records for registered players.
  \item Processes deposit and withdrawal operations on individual balances.
  \item Stores transaction records between players as an audit log.
  \item Enforces OAuth2 scope rules: \code{api:read} for all \code{GET} requests, \code{api:write} for all mutating requests.
  \item Returns standardized RFC 9457 \code{ProblemDetail} payloads for all error conditions.
  \item Exposes interactive API documentation via Swagger UI on the local profile.
\end{itemize}

% =============================================================================
\section{Architecture and Design}
% =============================================================================

The codebase follows a classic layered architecture. Each layer communicates only with the layer directly below it; no layer skips levels or introduces circular dependencies.

\subsection{Layer Descriptions}

\textbf{Controller layer} (\code{controller/}): Exposes REST endpoints under \code{/api/**}. Validates request DTOs using Bean Validation annotations and returns typed DTO responses with appropriate HTTP status codes. Contains no business logic.

\textbf{Service layer} (\code{service/}): Encapsulates business rules and transactional behavior. Throws typed domain exceptions (e.g., insufficient funds, duplicate player) that the exception handler maps to \code{ProblemDetail} responses.

\textbf{Repository layer} (\code{repository/}): Spring Data JPA repositories with custom queries for top-balance ranking and pessimistic locking on concurrent balance updates.

\textbf{Domain/model layer} (\code{model/}): JPA entities — \code{Player}, \code{Balance}, \code{Transaction}. Contains no behavior beyond JPA annotations.

\textbf{Cross-cutting concerns}: \code{config/} holds the security and OpenAPI configuration; \code{exceptions/} holds the centralized exception handler; \code{mapper/} holds entity-to-DTO mapping components.

\subsection{Security Model}

The API is a stateless OAuth2 resource server. JWTs are validated locally against the issuer URI configured in \code{AUTH\_ISSUER\_URI}. No session is created.

\rowcolors{2}{CraftLight}{white}
\begin{tabularx}{\linewidth}{@{} L{3.2cm} X @{}}
  \toprule
  {\tableheadstyle Scope} & {\tableheadstyle Permitted operations} \\
  \midrule
  \code{SCOPE\_api:read}  & \code{GET /api/**} \\
  \code{SCOPE\_api:write} & \code{POST}, \code{PUT}, \code{PATCH}, \code{DELETE} on \code{/api/**} \\
  \midrule
  \multicolumn{2}{@{}l}{\small\textit{Public (no token required): \code{GET /actuator/health}, \code{/swagger-ui/**}, \code{/v3/api-docs/**}}} \\
  \bottomrule
\end{tabularx}

\subsection{Error Contract}

All errors are returned as RFC 9457 \code{ProblemDetail} with the following extension fields:

\rowcolors{2}{CraftLight}{white}
\begin{tabularx}{\linewidth}{@{} L{2.8cm} X @{}}
  \toprule
  {\tableheadstyle Field} & {\tableheadstyle Description} \\
  \midrule
  \code{type}      & URI identifying the error category: \code{.../validation}, \code{.../business-rule}, or \code{.../internal}. \\
  \code{detail}    & Human-readable explanation of the specific error. \\
  \code{status}    & HTTP status code as an integer. \\
  \code{timestamp} & ISO 8601 timestamp of when the error occurred. \\
  \code{path}      & The request path that produced the error. \\
  \code{errors}    & Field-level validation map (present on \code{.../validation} errors only). \\
  \bottomrule
\end{tabularx}

\subsection{API Endpoints}

\rowcolors{2}{CraftLight}{white}
\begin{tabularx}{\linewidth}{@{} L{1.2cm} L{6.0cm} X @{}}
  \toprule
  {\tableheadstyle Method} & {\tableheadstyle Path} & {\tableheadstyle Description} \\
  \midrule
  \multicolumn{3}{@{}l}{\small\textit{\textbf{Players}}} \\
  GET    & \code{/api/players}                   & List all players. \\
  GET    & \code{/api/players/\{uuid\}}           & Get player by UUID. \\
  GET    & \code{/api/players/name/\{name\}}      & Get player by display name. \\
  POST   & \code{/api/players}                   & Register a new player. \\
  \midrule
  \multicolumn{3}{@{}l}{\small\textit{\textbf{Balances}}} \\
  GET    & \code{/api/balances}                  & List all balances. \\
  GET    & \code{/api/balances/\{uuid\}}          & Get balance by player UUID. \\
  GET    & \code{/api/balances/top}              & Top balances (\code{?limit=}, clamped 1–20, default 10). \\
  POST   & \code{/api/balances}                  & Create a balance record. \\
  PUT    & \code{/api/balances/\{uuid\}/set}      & Overwrite a balance amount. \\
  POST   & \code{/api/balances/\{uuid\}/deposit}  & Deposit funds. \\
  POST   & \code{/api/balances/\{uuid\}/withdraw} & Withdraw funds. \\
  \midrule
  \multicolumn{3}{@{}l}{\small\textit{\textbf{Transactions}}} \\
  GET    & \code{/api/transactions}              & List all transactions. \\
  GET    & \code{/api/transactions/id/\{id\}}     & Get transaction by ID. \\
  GET    & \code{/api/transactions/from/\{uuid\}} & List outgoing transactions for a player. \\
  GET    & \code{/api/transactions/to/\{uuid\}}   & List incoming transactions for a player. \\
  POST   & \code{/api/transactions}              & Store a transaction record. \\
  \bottomrule
\end{tabularx}

\important{\code{POST /api/transactions} stores a record only. It does \textbf{not} call balance transfer logic. Balance state and the transaction ledger are not updated atomically. See Known Limitations.}

% =============================================================================
\section{Tech Stack}
% =============================================================================

\rowcolors{2}{CraftLight}{white}
\begin{tabularx}{\linewidth}{@{} L{3.0cm} L{5.2cm} X @{}}
  \toprule
  {\tableheadstyle Category} &
  {\tableheadstyle Technology} &
  {\tableheadstyle Notes} \\
  \midrule
  Language      & Java 17                         & \\
  Framework     & Spring Boot 3.5                 & Web, Validation, Actuator \\
  Security      & Spring Security, OAuth2 Resource Server & Stateless JWT validation \\
  Persistence   & Spring Data JPA                 & Hibernate provider \\
  Migrations    & Flyway                          & Docker profile only \\
  Database      & PostgreSQL (Docker profile)     & \\
  Database      & H2 in-memory (local profile)    & Schema recreated on each startup \\
  API Docs      & springdoc-openapi               & Swagger UI at \code{/api-docs} \\
  Testing       & JUnit 5, Mockito, Spring Test   & \\
  Build Tool    & Gradle Wrapper                  & \\
  Packaging     & Docker multi-stage image        & \\
  \bottomrule
\end{tabularx}

% =============================================================================
\section{Configuration}
% =============================================================================

\rowcolors{2}{CraftLight}{white}
\begin{tabularx}{\linewidth}{@{} L{4.8cm} L{2.0cm} C{1.6cm} X @{}}
  \toprule
  {\tableheadstyle Variable} &
  {\tableheadstyle Default} &
  {\tableheadstyle Required} &
  {\tableheadstyle Description} \\
  \midrule
  \code{SPRING\_PROFILES\_ACTIVE}     & ---                          & \textbf{Yes} & Set to \code{local} (H2) or \code{docker} (PostgreSQL). \\
  \code{AUTH\_ISSUER\_URI}            & \code{http://localhost:9000} & No           & JWT issuer URI. Must match the Authorization Server's configured issuer. \\
  \code{SPRING\_DATASOURCE\_URL}      & ---                          & Docker only  & JDBC connection string. \\
  \code{SPRING\_DATASOURCE\_USERNAME} & ---                          & Docker only  & Database username. \\
  \code{SPRING\_DATASOURCE\_PASSWORD} & ---                          & Docker only  & Database password. \\
  \bottomrule
\end{tabularx}

% =============================================================================
\section{Known Limitations}
% =============================================================================

\begin{itemize}
  \item \code{POST /api/transactions} does not atomically update balances. The ledger and balance state can diverge if the system fails between a balance operation and the subsequent transaction record write.
  \item List endpoints return all records in a single response with no pagination or filtering support. This will degrade under large datasets.
  \item Integration tests run against H2 in-memory, not a real PostgreSQL instance. Schema and query compatibility gaps may exist.
  \item No CI pipeline is configured.
\end{itemize}

% =============================================================================
\section{Roadmap}
% =============================================================================

\begin{itemize}
  \item Integrate \code{POST /api/transactions} with \code{BalanceService.transfer()} so ledger and balance updates occur atomically within a single database transaction.
  \item Add cursor-based or offset pagination and server-side filtering to all list endpoints.
  \item Add integration tests against real PostgreSQL with actual JWT-bearing requests.
  \item Configure a CI pipeline with lint, typecheck, build, and test stages.
\end{itemize}
